

\chapter{数理逻辑}

\section{逻辑学简介}

\textbf{思维}：是人类对客观事物的反映，是人类特有的心理活动。
\begin{enumerate}
    \item \textbf{思维的内容}：思维所反映的具体对象、对象的属性、对象之间的关系，也就是思维的具体素材、具体含义
    \item \textbf{思维的形式}：思维内容各组成部分之间的联结方式、组织结构，是去掉具体内容之后，所有同类思维都共同具备的、通用的、稳定的框架
\end{enumerate}
\vspace{1em}

\begin{note}
    思维的存在本身是一个哲学问题，这里不展开讨论，假设人类思维是客观存在的，且所有人类思维的基本形式都是相同的，思维的差异仅体现在内容上。
\end{note}

\textbf{语言 Language} 是思维物化的结果，用于交换和记录思维。语言是一个符号系统，包含三个要素：
\begin{enumerate}
    \item \textbf{基本符号 Basic Symbol}：构成语言的最小基本单位。比如：英文的 26 个字母和标点符号，中文的所有汉字和标点符号等；
    \item \textbf{语法规则 Syntactic Rules}：规定哪些由基本符号排列而成的字符串是合法的词、词组和语句，又可分为词法规则与语句规则。比如：英文中 “good” 是一个合法的词，符合词法规则，而 “oogd” 则不是；再比如，“The book is good.” 是一个合法的句子，符合句法规则，而 “good The is book” 则不是；
    \item \textbf{语义规则 Semantic Rules}：规定合法的词、词组和语句的意义。比如：英文中 “good” 的意义是“好的”，而 “bad” 的意义是“不好的”；再比如，“The book is good.” 的意义是“这本书是好的”，而 “The book is bad.” 的意义则是“这本书是不好的”。
\end{enumerate}
语言是一种人造的产物，又可以分为：
\begin{enumerate}
    \item \textbf{自然语言 Natural Language}：是人类社会在长期交往中自然形成、协同演化、代代相传的符号系统，用于日常思维表达，如汉语、英语和日语等。自然语言因为“语境”的不同，通常会产生歧义；语境不是语言的构成要素，而是理解自然语言的背景；
    \item \textbf{人工语言 Artificial Language}：基于明确的设计目标和使用场景，严格规定全部符号、语法和语义。不存在自然演化的过程，理解人工语言也不需要考虑语境，理论上不存在歧义，比如数学的符号语言、计算机程序语言等。
\end{enumerate}
\vspace{1em}
\begin{note}
    \textbf{语言学 Linguistics}是专门研究人类语言的学科，关于思维和语言的关系学界并无定论，这里只采取了主流语言学家的观点，即思维可以不依赖语言而独立存在，而语言是思维的外化，用于交流和记录的工具。
\end{note}

\textbf{逻辑学 Logic} 关注理性思维，研究思维的形式，也即研究命题的性质及其推理有效性。

\begin{note}
    逻辑学和语言的关系分为两个层次，一是逻辑学研究的对象（命题及其推理）的语言，称为对象语言，比如命题逻辑中使用的人工语言 $\mathcal{L}_0$；逻辑学讨论研究对象的语言，称为元语言，比如我们在这里使用的自然语言。
\end{note}

\subsection{词项}

\subsection{命题}

\subsection{推理}



\newpage